% Working manuscript: no current-study results have been supplied.
% Keep the presentation free of subsections, per the repository guidelines.
% Background is integrated into the presentation rather than a separate Introduction.
\begin{PRSabstract}
Transoral robotic surgery offers a platform for working within the confined oral cavity, but instrument size remains a practical concern for cleft palate repair. Miniature 3 mm wristed instruments previously evaluated in a synthetic cleft palate simulator provide the basis for a live porcine investigation of simulated repair. The study premise describes one surgeon and eight pigs, with four animals assigned to each instrument size: the laboratory's 3 mm tools and standard 8 mm da Vinci tools. The planned evaluation combines stereo endoscopic video and robot recordings to assess completion of the repair, immediate closure quality, instrument obstruction of the surgical view, and suture duration. Tissue-force comparisons depend on whether logged effort can support a validated estimate. This investigation addresses operative performance in living tissue and the practical value of instrument miniaturization for transoral reconstruction. [Insert verified completion counts, closure findings, comparative estimates, and a conclusion supported by those results after analysis.]
\PRSkeywords{transoral robotic surgery; simulated cleft palate repair; miniature instruments; porcine model}
\end{PRSabstract}
\PRSstartbody

\section{Miniature Instruments for Simulated Cleft Palate Repair}
Cleft palate repair requires tissue dissection and suturing in a restricted oral workspace. Robotic visualization and wrist articulation may assist these tasks, although large instruments can obstruct the view and contact adjacent structures. Maguire and colleagues developed compact, cable-driven, pin-jointed 3 mm instruments for the da Vinci Research Kit and evaluated them using a high-fidelity synthetic cleft palate simulator. Their study combined geometric reachability analysis with physical experiments and reported fewer collisions and reduced image obstruction compared with 8 mm instruments.\supercite{maguire2025}

The present innovation applies these miniature instruments to simulated cleft palate repair in a live porcine model. The shared tool concept motivates this study, which examines completion and quality of repair in living tissue together with practical measures of visualization and suturing performance. The study premise specifies one surgeon and eight pigs, four per instrument size. [Confirm animal characteristics, allocation and trial order, number of repairs per animal, robotic platform for each group, and prior operator experience.]

[Describe animal preparation, oral exposure, creation of the simulated defect, instrument configuration, dissection, and documented nasal, muscle, and oral closure steps. Specify suture material, needle, tool exchanges, and any planned manual steps. Add animal ethics approval and perioperative details once verified.]

Stereo endoscopic recordings and robot data stored in rosbags will be matched to each repair and synchronized before annotation. Feasibility will be assessed by completion of the intended steps, unplanned conversion to larger instruments or manual surgery, instrument failures, dropped needles, and setup time. Immediate closure assessment will address visible gaps, suture spacing, and bite depth where measurement is supported. A clinically justified comparability margin and scoring rules must be specified before outcome review.

Instrument obstruction will be expressed as the fraction of a defined endoscopic image region occupied by instruments during comparable tasks. Suture duration will run from first needle--tissue contact to knot tightening or suture cutting, using a consistent endpoint and separate summaries for each closure layer. Supporting measures include observed collisions, clutching and repositioning, posterior access, grasps and regrasps, and operative-step duration. Robot motion may support path-length, smoothness, and tool-relative workspace summaries after coordinate and sampling checks.

[Insert observed completion denominators, closure assessments, and per-repair comparative summaries with uncertainty. Report all computed outcomes, including those favoring 8 mm instruments, and identify unavailable measurements.]

\section{Discussion}
The clinical motivation is to make dexterous transoral reconstruction practical in a small workspace while retaining sufficient capability for dissection and closure. The earlier simulator investigation establishes the engineering rationale; this porcine study addresses a subsequent preclinical question about operative performance in living tissue.\supercite{maguire2025} Successful completion, if demonstrated, must accompany any proposed advantage in visualization, speed, or tissue loading.

Image-area obstruction offers a practical view of instrument interference with the operative field. Its interpretation requires comparable camera position, magnification, task, and image region. Camera calibration is unavailable, so physical bite depth and spacing require an independent scale or measurement method. Absolute three-dimensional tissue reconstruction cannot be assumed from stereo recordings alone.

Logged joint effort is not a direct measurement of tissue force. Any force estimate requires verification of topic contents, units, instrument-specific transmission, friction, and calibration. Jaw effort describes a different interaction from needle-driving or traction force. Without validation, effort should be reported as a proxy; video deformation alone cannot establish lower force or reduced tissue injury.

This is a small, single-surgeon series. Repeated sutures within a repair must not be treated as independent animals, and a surgeon random effect cannot be estimated from one operator. Trial order, learning, anatomy, and setup may influence comparisons. Similar closure scores alone would not establish equivalence. Porcine simulated defects also limit inference about infant anatomy, healing, speech, and clinical safety. The final interpretation should therefore emphasize demonstrated technical feasibility and measured operative differences, with broader clinical claims reserved for further validation.
